\documentclass[11pt]{article}

\usepackage[a4paper,margin=2.5cm]{geometry}
\usepackage{amsmath,amssymb,amsthm}
\usepackage{hyperref}
\usepackage{microtype}

\title{Die Fibonacci-Zetafunktion als Summe geometrischer Resonatoren}
\author{Thomas Hoffbauer}
\date{\today}

\theoremstyle{definition}
\newtheorem{definition}{Definition}
\newtheorem{remark}{Bemerkung}

\begin{document}
	
	\maketitle
	
	\begin{abstract}
		Wir leiten eine Darstellung der Fibonacci-Zetafunktion als unendliche
		Summe geometrischer Resonatoren her.
		Ausgehend von der Binet-Formel entsteht eine Entwicklung,
		deren Pole durch einfache Resonanzbedingungen bestimmt werden.
		Die Darstellung macht sichtbar, dass die Fibonacci-Zeta
		als \"Uberlagerung elementarer geometrischer Reihen verstanden werden kann.
	\end{abstract}
	
	\section{Einleitung}
	
	Die klassische Riemannsche Zetafunktion
	\[
	\zeta(s)
	=
	\sum_{n\ge1}\frac1{n^s}
	\]
	ist urspr\"unglich f\"ur
	\[
	\operatorname{Re}(s)>1
	\]
	definiert.
	
	Eine nat\"urliche Variante entsteht,
	wenn man die nat\"urlichen Zahlen durch die Fibonacci-Zahlen ersetzt.
	
	Die Fibonacci-Zahlen wachsen exponentiell.
	Aus der Binet-Formel folgt
	\[
	F_n
	=
	\frac{\varphi^n-\psi^n}{\sqrt5},
	\]
	mit
	\[
	\varphi
	=
	\frac{1+\sqrt5}{2},
	\qquad
	\psi
	=
	\frac{1-\sqrt5}{2}.
	\]
	Daraus ergibt sich die Asymptotik
	\[
	F_n
	\sim
	\frac{\varphi^n}{\sqrt5}
	\qquad (n\to\infty).
	\]
	Somit gilt
	\[
	F_n^{-s}
	\sim
	5^{s/2}\varphi^{-ns}.
	\]
	Da die rechte Seite geometrisch abf\"allt,
	konvergiert die Reihe
	\[
	\sum_{n\ge1}F_n^{-s}
	\]
	absolut genau dann, wenn
	\[
	|\varphi^{-s}|<1.
	\]
	Dies ist \"aquivalent zu
	\[
	\operatorname{Re}(s)>0.
	\]
	
	\begin{definition}[Fibonacci-Zetafunktion]
		F\"ur komplexe Zahlen $s$ mit
		\[
		\operatorname{Re}(s)>0
		\]
		definieren wir
		\[
		\zeta_F(s)
		=
		\sum_{n\ge1}
		\frac1{F_n^s}.
		\]
	\end{definition}
	
	Unser Ziel ist eine Darstellung dieser Funktion
	als Summe elementarer Resonatoren.
	
	\section{Die Binet-Darstellung}
	
	Die geschlossene Formel von Binet lautet
	\[
	F_n
	=
	\frac{\varphi^n-\psi^n}{\sqrt5}.
	\]
	Mit
	\[
	\psi=-\varphi^{-1}
	\]
	erh\"alt man
	\[
	F_n
	=
	\frac{\varphi^n}{\sqrt5}
	\left(
	1-(-1)^n\varphi^{-2n}
	\right).
	\]
	Daher
	\[
	F_n^{-s}
	=
	5^{s/2}
	\varphi^{-ns}
	\left(
	1-(-1)^n\varphi^{-2n}
	\right)^{-s}.
	\]
	
	\section{Binomialentwicklung}
	
	\begin{definition}[Pochhammer-Symbol]
		F\"ur $m\ge0$ definiert man
		\[
		(s)_m
		=
		s(s+1)(s+2)\cdots(s+m-1),
		\]
		mit der Konvention
		\[
		(s)_0=1.
		\]
		Das Symbol wird auch als aufsteigende Fakult\"at
		(rising factorial) bezeichnet.
	\end{definition}
	
	F\"ur
	\[
	|u|<1
	\]
	gilt die verallgemeinerte Binomialreihe
	\[
	(1-u)^{-s}
	=
	\sum_{m\ge0}
	\frac{(s)_m}{m!}u^m.
	\]
	Setzt man
	\[
	u=(-1)^n\varphi^{-2n},
	\]
	so folgt
	\[
	F_n^{-s}
	=
	5^{s/2}
	\sum_{m\ge0}
	\frac{(s)_m}{m!}
	(-1)^{mn}
	\varphi^{-n(s+2m)}.
	\]
	
	\section{Vertauschung der Summationen}
	
	Durch Summation \"uber alle Fibonacci-Indizes erh\"alt man
	\[
	\zeta_F(s)
	=
	\sum_{n\ge1}F_n^{-s}.
	\]
	Einsetzen der Entwicklung liefert
	\[
	\zeta_F(s)
	=
	5^{s/2}
	\sum_{n\ge1}
	\sum_{m\ge0}
	\frac{(s)_m}{m!}
	(-1)^{mn}
	\varphi^{-n(s+2m)}.
	\]
	F\"ur
	\[
	\operatorname{Re}(s)>0
	\]
	konvergiert die entstehende Doppelsumme absolut.
	Daher ist nach dem Satz von Fubini bzw.\ Tonelli
	die Vertauschung der Summationsreihenfolge zul\"assig.
	
	Somit
	\[
	\zeta_F(s)
	=
	5^{s/2}
	\sum_{m\ge0}
	\frac{(s)_m}{m!}
	\sum_{n\ge1}
	\Bigl(
	(-1)^m
	\varphi^{-(s+2m)}
	\Bigr)^n.
	\]
	
	\section{Geometrische Resonatoren}
	
	Die innere Summe ist eine geometrische Reihe.
	
	F\"ur
	\[
	|q|<1
	\]
	gilt
	\[
	\sum_{n\ge1}q^n
	=
	\frac{q}{1-q}.
	\]
	Mit
	\[
	q
	=
	(-1)^m
	\varphi^{-(s+2m)}
	\]
	erh\"alt man
	\[
	\boxed{
		\zeta_F(s)
		=
		5^{s/2}
		\sum_{m\ge0}
		\frac{(s)_m}{m!}
		\frac{
			(-1)^m\varphi^{-(s+2m)}
		}{
			1-(-1)^m\varphi^{-(s+2m)}
		}.
	}
	\]
	
	Dies ist die Resonatordarstellung der Fibonacci-Zetafunktion.
	
	\section{Resonanzbedingungen}
	
	Definiere
	\[
	R_m(s)
	=
	\frac{
		(-1)^m\varphi^{-(s+2m)}
	}{
		1-(-1)^m\varphi^{-(s+2m)}
	}.
	\]
	Jeder Term besitzt die Form eines elementaren Resonators.
	
	Pole entstehen genau dann, wenn
	\[
	1-(-1)^m\varphi^{-(s+2m)}
	=0.
	\]
	Dies ist \"aquivalent zu
	\[
	\varphi^{s+2m}
	=
	(-1)^m.
	\]
	Die Fibonacci-Zetafunktion erscheint somit als \"Uberlagerung
	einer unendlichen Familie geometrischer Resonatoren,
	deren Resonanzlagen durch den Goldenen Schnitt bestimmt werden.
	
	\section{Schlussbemerkung}
	
	Die erhaltene Darstellung zeigt,
	dass die Fibonacci-Zetafunktion als Summe einfacher
	geometrischer Resonatoren interpretiert werden kann.
	
	Der Goldene Schnitt
	\[
	\varphi=\frac{1+\sqrt5}{2}
	\]
	\"ubernimmt dabei die Rolle einer fundamentalen Skala.
	
	Weitere Untersuchungen k\"onnten die analytische Fortsetzung,
	die genaue Struktur der Resonanzpole sowie Verallgemeinerungen
	auf andere rekursiv definierte Zahlenfolgen behandeln.
	
\end{document}
